\chapter{Integral Riemann dan Teorema Dasar Kalkulus}

Sebuah program menghitung luas di bawah $y=x^2$ pada $[0,1]$ dengan 1000 potong sama
lebar. Bila tinggi persegi panjang diambil di ujung kiri, hasilnya $0.3328335$; bila di
ujung kanan, $0.3338335$. Selisihnya $0.001$, dan tidak ada alasan memilih salah satu.
Nilai yang benar adalah $1/3$, yang tidak muncul pada kedua hitungan itu. Menambah
potongan menjadi sejuta tidak menolong: kedua jawaban tetap mengapit, tidak pernah
menangkap.

Yang diperlukan bukan partisi lebih halus, melainkan cara membaca dua hampiran itu
sekaligus. Kalau selisihnya dapat dibuat lebih kecil dari bilangan positif mana pun,
maka ada satu bilangan yang terkurung di antaranya, dan itulah integralnya. Seluruh
bab ini berdiri di atas gagasan sederhana itu.

Bab 12 mengajari kita membongkar perubahan menjadi laju: dari posisi ke kecepatan,
dari fungsi ke turunan. Yang belum ada adalah jalan pulang. Kita tahu kecepatan pada
detik ke-$t$ adalah $v(t)$, tetapi belum tahu berapa jarak yang ditempuh selama satu
menit. Intuisinya jelas: jarak adalah jumlah kecepatan dikali potongan waktu kecil.
Yang belum jelas adalah cara menjumlahkan tak berhingga banyak suku itu dengan benar.
Bab ini menjawabnya, dan jawabannya mengejutkan: jalan pulang itu adalah antiturunan.

\begin{untukanak}
Bayangkan kita ingin tahu luas sebuah danau yang bentuknya tidak beraturan. Tidak ada
rumus untuk bentuk seperti itu. Yang kita punya hanya kertas berpetak bening. Kita
letakkan kertas itu di atas gambar danau, lalu hitung berapa petak yang seluruhnya
berada di dalam air. Hasilnya terlalu kecil. Kita hitung lagi berapa petak yang
tersentuh air, walau hanya sedikit. Hasilnya terlalu besar. Luas sebenarnya ada di
antara dua angka itu.

Sekarang potong petaknya menjadi sepuluh kali lebih kecil. Angka bawah naik, angka atas
turun, dan keduanya makin rapat. Kalau petaknya sangat kecil, kedua angka hampir sama,
dan kita bisa menyebut luas danau sampai ketelitian berapa pun.

Ada satu syarat yang mudah terlewat. Kalau danau punya pulau karang sangat tipis yang
tersebar di mana-mana, petak yang menyentuhnya bisa membuat taksiran atas tetap jauh
dari taksiran bawah. Matematikawan menyebut syarat itu keterintegralan. Di bab ini kita
akan tahu mengapa dua taksiran yang mengapit bisa dipercaya, dan kapan taksiran itu
gagal.
\end{untukanak}

\section{Partisi, jumlah atas, dan jumlah bawah}
\tingkat{dasar}

Kita ingin memberi arti pada luas di bawah grafik tanpa mengandaikan luas itu sudah
punya arti: tutup daerahnya dengan persegi panjang dari dalam dan dari luar, lalu
persempit selisihnya.

\begin{definisi}[Partisi dan jumlah atas-bawah]\index{partisi}
Misalkan $[a,b]$ selang tertutup. Sebuah \emph{partisi} $\mathcal{P}$ dari $[a,b]$
adalah himpunan berhingga titik
\[
a=x_0<x_1<\cdots<x_n=b .
\]
Lebar potongan ke-$i$ adalah $\Delta x_i=x_i-x_{i-1}$, dan \emph{mesh} partisi itu
adalah lebar potongan terlebar, $\|\mathcal{P}\|=\max_i\Delta x_i$.

Misalkan $f$ terbatas pada $[a,b]$. Pada potongan ke-$i$ tulis
\[
M_i=\sup\{f(x):x_{i-1}\le x\le x_i\},\qquad
m_i=\inf\{f(x):x_{i-1}\le x\le x_i\}.
\]
\emph{Jumlah atas} dan \emph{jumlah bawah} partisi $\mathcal{P}$ adalah
\begin{equation}
U(\mathcal{P},f)=\sum_{i=1}^{n}M_i\,\Delta x_i,
\qquad
L(\mathcal{P},f)=\sum_{i=1}^{n}m_i\,\Delta x_i .
\end{equation}
\end{definisi}

Kita memakai supremum dan infimum, bukan nilai terbesar dan terkecil, karena $f$ belum
tentu mencapai ekstremnya pada potongan: fungsi yang bernilai $x$ untuk $x<1$ dan $0$
di $x=1$ terbatas pada $[0,1]$, tetapi nilai terbesarnya pada potongan terakhir tidak
ada, sedangkan supremumnya ada, yaitu $1$.

Jumlah atas selalu paling sedikit sama dengan jumlah bawah, sebab $M_i\ge m_i$ untuk
setiap $i$. Untuk $f(x)=x^2$ pada $[0,1]$ dengan empat potongan sama lebar, jumlah
bawahnya
\[
L=\frac{1}{4}\Bigl(0+\frac{1}{16}+\frac{4}{16}+\frac{9}{16}\Bigr)=\frac{7}{32},
\]
dan jumlah atasnya $\frac{15}{32}$. Keduanya mengapit nilai yang dicari, dengan selisih
$\frac14$ yang masih besar. Bila $f$ bernilai negatif, jumlah ini menghitung luas
bertanda, dan tanda minusnya bagian dari definisi, bukan kekeliruan.

Gambar~\ref{fig:rect-halus} memperlihatkan dua nilai $n$ sekaligus, dengan ujung kanan
sebagai tinggi. Pada $n=4$ tangganya kasar dan $U_4=0.46875$; pada $n=16$ geriginya
halus dan $U_{16}=0.3652343750$. Jarak tangga ke garis $1/3$ menyusut dari $0.1354$
menjadi $0.0319$.

\begin{figure}[htbp]
\centering
\begin{tikzpicture}
\begin{axis}[width=0.47\linewidth,height=4.6cm,
  axis lines=left,xmin=0,xmax=1.02,ymin=0,ymax=1.12,
  xtick={0,0.5,1},ytick={0,0.5,1},
  tick label style={font=\scriptsize},label style={font=\scriptsize},
  xlabel={$x$},
  title={\scriptsize $n=4$: $U_4=0.46875$}]
\pgfplotsinvokeforeach{1,2,3,4}{
  \draw[fill=layerteal!18,draw=layerteal!65,line width=0.4pt]
    (axis cs:{(#1-1)/4},0) rectangle (axis cs:{#1/4},{#1*#1/16});
}
\draw[densely dashed,layergrey,line width=0.6pt]
  (axis cs:0,0.33333) -- (axis cs:1,0.33333);
\addplot[thick,accent,domain=0:1,samples=80] {x^2};
\node[font=\scriptsize,text=layergrey,anchor=west] at (axis cs:0.06,0.40) {$1/3$};
\end{axis}
\end{tikzpicture}\hfill
\begin{tikzpicture}
\begin{axis}[width=0.47\linewidth,height=4.6cm,
  axis lines=left,xmin=0,xmax=1.02,ymin=0,ymax=1.12,
  xtick={0,0.5,1},ytick={0,0.5,1},
  tick label style={font=\scriptsize},label style={font=\scriptsize},
  xlabel={$x$},
  title={\scriptsize $n=16$: $U_{16}=0.3652343750$}]
\pgfplotsinvokeforeach{1,...,16}{
  \draw[fill=layerteal!18,draw=layerteal!65,line width=0.3pt]
    (axis cs:{(#1-1)/16},0) rectangle (axis cs:{#1/16},{#1*#1/256});
}
\draw[densely dashed,layergrey,line width=0.6pt]
  (axis cs:0,0.33333) -- (axis cs:1,0.33333);
\addplot[thick,accent,domain=0:1,samples=80] {x^2};
\node[font=\scriptsize,text=layergrey,anchor=west] at (axis cs:0.06,0.40) {$1/3$};
\end{axis}
\end{tikzpicture}
\caption{Penghampiran $f(x)=x^{2}$ pada $[0,1]$ dengan persegi panjang berujung kanan.
Pada $n=4$ jumlah atas $U_4=0.46875$, pada $n=16$ jumlah atas $U_{16}=0.3652343750$,
dan keduanya masih di atas nilai sebenarnya $1/3=0.333\ldots$ (garis putus-putus).
Selisih $U_n-L_n$ tepat $1/n$, jadi naik dari $n=4$ ke $n=16$ memotong jarak tangga ke
garis itu menjadi seperempatnya.}
\label{fig:rect-halus}
\end{figure}


Ada dua kenyataan yang akan dipakai berulang kali. Yang pertama: memperhalus partisi
tidak memperburuk taksiran.

\begin{lemma}[Penghalusan]\index{penghalusan}
Jika $\mathcal{P}^{*}$ adalah penghalusan dari $\mathcal{P}$, artinya
$\mathcal{P}\subseteq\mathcal{P}^{*}$, maka
\[
L(\mathcal{P},f)\le L(\mathcal{P}^{*},f)\le U(\mathcal{P}^{*},f)\le U(\mathcal{P},f).
\]
\end{lemma}

\begin{proof}
Cukup dibuktikan untuk penghalusan yang menambah satu titik $c$ di dalam potongan
$[x_{i-1},x_i]$, sebab penghalusan umum diperoleh dengan mengulang langkah itu.
Misalkan $m'$ dan $m''$ infimum $f$ pada $[x_{i-1},c]$ dan pada $[c,x_i]$. Kedua
selang itu bagian dari $[x_{i-1},x_i]$, sehingga $m'\ge m_i$ dan $m''\ge m_i$.
Akibatnya
\[
m'(c-x_{i-1})+m''(x_i-c)\ \ge\ m_i(c-x_{i-1})+m_i(x_i-c)=m_i\,\Delta x_i ,
\]
jadi jumlah bawah tidak turun. Argumen yang sama dengan supremum memberi $M'\le M_i$
dan $M''\le M_i$, sehingga jumlah atas tidak naik.
\end{proof}

Kenyataan kedua lebih kuat: dua partisi berbeda tetap dapat dibandingkan, asal yang
dibandingkan adalah jumlah bawah dengan jumlah atas.

\begin{lemma}
Untuk sebarang dua partisi $\mathcal{P}_1$ dan $\mathcal{P}_2$ dari $[a,b]$ berlaku
\[
L(\mathcal{P}_1,f)\le U(\mathcal{P}_2,f).
\]
\end{lemma}

\begin{proof}
Ambil $\mathcal{P}$ penghalusan bersama, yaitu gabungan titik kedua partisi. Menurut
lemma sebelumnya, $L(\mathcal{P}_1,f)\le L(\mathcal{P},f)$ dan
$U(\mathcal{P},f)\le U(\mathcal{P}_2,f)$. Karena $L(\mathcal{P},f)\le U(\mathcal{P},f)$
selalu benar, rantai ketaksamaan itu memberi hasil yang diminta.
\end{proof}

Lemma itu memastikan setiap jumlah bawah merupakan batas bawah bagi seluruh jumlah
atas. Karena itu himpunan semua jumlah bawah terbatas di atas, dan semua jumlah atas
terbatas di bawah; kedua bilangan berikut selalu ada.

\begin{definisi}[Terintegralkan Riemann]\index{terintegralkan riemann}
Untuk $f$ terbatas pada $[a,b]$, \emph{integral bawah} dan \emph{integral atas}
didefinisikan sebagai
\[
\underline{I}(f)=\sup_{\mathcal{P}}L(\mathcal{P},f),
\qquad
\overline{I}(f)=\inf_{\mathcal{P}}U(\mathcal{P},f),
\]
dengan supremum dan infimum diambil atas seluruh partisi $[a,b]$. Karena
$\underline{I}(f)\le\overline{I}(f)$, fungsi $f$ disebut \emph{terintegralkan Riemann}
pada $[a,b]$ bila kedua bilangan itu sama. Nilai bersamanya ditulis
$\int_a^b f$ atau $\int_a^b f(x)\,dx$.
\end{definisi}

Perhatikan apa yang tidak ada pada definisi itu: tidak ada titik sampel dan tidak ada
limit atas $n$. Ukuran keberhasilan sebuah hampiran adalah selisih kedua kurungannya,
bukan kedekatannya pada tebakan.

\kamus{partisi}{partition}{Cara memotong selang menjadi beberapa potongan, seperti
membagi sebatang cokelat dengan pisau, bukan mengukur berapa panjang setiap potongan
dengan tepat.}

\kamus{jumlah atas}{upper sum}{Taksiran yang sengaja dibesar-besarkan, seperti
menimbang barang bersama kotaknya.}

\begin{bacakeras}
Dua suara. Sebut yang pertama \emph{Pemotong}, yang kedua \emph{Penjumlah}.

Pemotong: Aku punya sebatang pita. Aku potong jadi empat.
Penjumlah: Aku jumlah keempat potong itu. Hasilnya terlalu kecil.
Pemotong: Aku potong lagi. Jadi delapan.
Penjumlah: Hasilnya naik sedikit. Masih terlalu kecil.
Pemotong: Enam belas.
Penjumlah: Naik lagi. Makin dekat, belum sampai.
Pemotong: Seratus.
Penjumlah: Dekat sekali.
Pemotong: Seribu.
Penjumlah: Lebih dekat lagi, tetap belum sampai.

Pemotong: Kapan sampai?
Penjumlah: Tidak pernah, kalau kau hanya menambah potongan. Berhenti mengejar angka,
mulai mengejar kurungan.
Pemotong: Kurungan apa?
Penjumlah: Potongan yang kurus memberi batas bawah. Potongan yang gemuk memberi batas
atas. Kalau dua batas itu bertemu di satu titik, titik itulah jawabannya.
Pemotong: Jadi potonganku tidak pernah cukup?
Penjumlah: Tidak pernah. Potongan hanya membuat dua batas itu berdekatan. Yang menjawab
bukan potongan terakhir, melainkan pertemuan dua batas.

Sekarang ulangi satu baris ini: potong tipis, kurung rapat, batas bertemu.
\end{bacakeras}

\section{Kriteria keterintegralan}
\tingkat{menengah}

Definisi tadi menyebut supremum dan infimum atas seluruh partisi, dan partisi itu tak
berhingga banyaknya. Kriteria berikut menggantinya dengan satu pertidaksamaan yang
dapat diperiksa pada satu partisi saja.

\begin{teorema}[Kriteria keterintegralan]\index{kriteria keterintegralan}
Fungsi terbatas $f$ terintegralkan Riemann pada $[a,b]$ jika dan hanya jika untuk
setiap $\varepsilon>0$ terdapat partisi $\mathcal{P}$ dengan
\[
U(\mathcal{P},f)-L(\mathcal{P},f)<\varepsilon .
\]
\end{teorema}

\begin{proof}
Untuk setiap partisi berlaku
$L(\mathcal{P},f)\le\underline{I}(f)\le\overline{I}(f)\le U(\mathcal{P},f)$; ketaksamaan
pertama dari definisi supremum, yang terakhir dari definisi infimum.

$(\Rightarrow)$ Andaikan $f$ terintegralkan dengan nilai $I$, dan ambil
$\varepsilon>0$. Menurut definisi supremum, ada $\mathcal{P}_1$ dengan
$L(\mathcal{P}_1,f)>I-\varepsilon/2$. Menurut definisi infimum, ada $\mathcal{P}_2$
dengan $U(\mathcal{P}_2,f)<I+\varepsilon/2$. Ambil $\mathcal{P}$ penghalusan bersama
keduanya. Lemma penghalusan memberi
\[
U(\mathcal{P},f)-L(\mathcal{P},f)
<\Bigl(I+\tfrac{\varepsilon}{2}\Bigr)-\Bigl(I-\tfrac{\varepsilon}{2}\Bigr)
=\varepsilon .
\]

$(\Leftarrow)$ Andaikan syarat itu dipenuhi. Untuk setiap $\varepsilon>0$ ada
$\mathcal{P}$ dengan $U(\mathcal{P},f)-L(\mathcal{P},f)<\varepsilon$, sehingga
\[
0\le\overline{I}(f)-\underline{I}(f)\le U(\mathcal{P},f)-L(\mathcal{P},f)<\varepsilon .
\]
Bilangan $\overline{I}(f)-\underline{I}(f)$ tidak negatif dan lebih kecil dari setiap
$\varepsilon>0$, jadi bilangan itu nol. Kedua integral sama, dan $f$ terintegralkan.
\end{proof}

Kriteria itu mengubah pertanyaan tentang tak berhingga banyak partisi menjadi satu
pertidaksamaan. Akibat langsungnya: bila selisihnya kecil pada satu partisi, selisih
itu tetap kecil pada setiap penghalusannya, sebab lemma penghalusan hanya dapat
mempersempitnya.

Sebelum mencari kelas fungsi yang lolos kriteria, kita catat sifat integral yang akan
dipakai terus. Semuanya turunan dari definisi jumlah atas dan bawah, bukan asumsi
baru.

\begin{proposisi}[Sifat dasar integral]\index{integral, sifat dasar}
Misalkan $f$ dan $g$ terintegralkan Riemann pada $[a,b]$. Maka:
(a) $f+g$ terintegralkan dan $\int_a^b(f+g)=\int_a^b f+\int_a^b g$;
(b) jika $f\le g$ pada $[a,b]$ maka $\int_a^b f\le\int_a^b g$;
(c) jika $a<c<b$ maka $f$ terintegralkan pada $[a,c]$ dan pada $[c,b]$, dan
$\int_a^b f=\int_a^c f+\int_c^b f$;
(d) $|f|$ terintegralkan dan $\bigl|\int_a^b f\bigr|\le\int_a^b|f|$;
(e) untuk setiap konstanta $k$, $\int_a^b k\,dx=k(b-a)$.
\end{proposisi}

\begin{proof}
(a) Pada setiap potongan, $\sup(f+g)\le\sup f+\sup g$ dan $\inf(f+g)\ge\inf f+\inf g$.
Jadi $U(\mathcal{P},f+g)\le U(\mathcal{P},f)+U(\mathcal{P},g)$ dan
$L(\mathcal{P},f+g)\ge L(\mathcal{P},f)+L(\mathcal{P},g)$. Ambil $\varepsilon>0$ dan
pilih $\mathcal{P}$ yang membuat $U-L$ kurang dari $\varepsilon/2$ untuk $f$ dan $g$
sekaligus. Maka
\[
U(\mathcal{P},f+g)-L(\mathcal{P},f+g)
\le\bigl[U(\mathcal{P},f)-L(\mathcal{P},f)\bigr]
+\bigl[U(\mathcal{P},g)-L(\mathcal{P},g)\bigr]<\varepsilon ,
\]
sehingga $f+g$ terintegralkan. Nilainya diperoleh dengan mengurung dua arah, karena
$\varepsilon$ sebarang.

(b) Bila $f\le g$, maka $M_i(f)\le M_i(g)$ pada setiap potongan, sehingga
$U(\mathcal{P},f)\le U(\mathcal{P},g)$ untuk semua $\mathcal{P}$. Ambil infimum kedua
ruas.

(c) Bila $\mathcal{P}$ memuat $c$, jumlahnya terpecah menjadi bagian $[a,c]$ dan
$[c,b]$, dengan $L(\mathcal{P},f)=L(\mathcal{P}_1,f)+L(\mathcal{P}_2,f)$ dan hal sama
untuk $U$. Supremum memberi
$\underline{I}(f)\le\underline{I}_{[a,c]}(f)+\underline{I}_{[c,b]}(f)$, infimum
memberi arah sebaliknya. Untuk partisi tanpa $c$, tambahkan $c$ lebih dulu; lemma
penghalusan menjaga ketaksamaan.

(d) Karena $-|f|\le f\le|f|$, bagian (b) dan (a) memberi
$-\int|f|\le\int f\le\int|f|$. Keterintegralkan $|f|$ menyusul dari
$\bigl||f(x)|-|f(y)|\bigr|\le|f(x)-f(y)|$, yang memberi
$M_i(|f|)-m_i(|f|)\le M_i(f)-m_i(f)$.

(e) Pada setiap potongan, $M_i=m_i=k$, jadi $U=L=k(b-a)$ untuk semua partisi.
\end{proof}

Bagian (c) adalah alasan kita boleh memotong selang di titik mana pun; tanpanya fungsi
akumulasi tidak dapat didefinisikan dengan tenang.

\section{Kelas fungsi yang terintegralkan}
\tingkat{menengah}

Kriteria pada bagian sebelumnya adalah alat uji. Sekarang kita cari kelas fungsi yang
lolos uji itu tanpa harus mencari partisi khusus untuk setiap kasus. Dua kelas paling
berguna adalah fungsi kontinu dan fungsi monoton.

\begin{akibat}[Kekontinuan menjamin keterintegralan]
Jika $f$ kontinu pada $[a,b]$, maka $f$ terintegralkan Riemann pada $[a,b]$.
\end{akibat}

\begin{proof}
Ambil $\varepsilon>0$. Karena $[a,b]$ kompak, $f$ kontinu seragam di sana (Bab 11):
ada $\delta>0$ sehingga $|f(x)-f(t)|<\varepsilon/(b-a)$ bila $|x-t|<\delta$. Pilih
partisi dengan mesh kurang dari $\delta$. Pada setiap potongan, untuk semua $x,t$ di
potongan itu berlaku $|f(x)-f(t)|<\varepsilon/(b-a)$, sehingga
$M_i-m_i\le\varepsilon/(b-a)$. Maka
\[
U(\mathcal{P},f)-L(\mathcal{P},f)
=\sum_{i=1}^{n}(M_i-m_i)\,\Delta x_i
\le\frac{\varepsilon}{b-a}\sum_{i=1}^{n}\Delta x_i=\varepsilon .
\]
Kriteria keterintegralan menutup bukti.
\end{proof}

Perhatikan peran kekontinuan seragam: satu $\delta$ harus bekerja untuk semua potongan
sekaligus, sedangkan kekontinuan biasa hanya memberi $\delta$ yang bergantung pada
titik. Di sinilah Bab 11 dipakai, bukan sebagai hiasan.

Kelas kedua tidak memerlukan kekontinuan: fungsi monoton boleh melompat, asal
lompatannya berhingga banyak.

\begin{teorema}[Monoton menjamin keterintegralan]
Jika $f$ monoton pada $[a,b]$, maka $f$ terintegralkan Riemann pada $[a,b]$.
\end{teorema}

\begin{proof}
Andaikan $f$ naik. Bagi $[a,b]$ menjadi $n$ potongan sama lebar, $\Delta x=(b-a)/n$.
Pada potongan ke-$i$, nilai terkecil ada di ujung kiri dan nilai terbesar di ujung
kanan, sehingga $M_i-m_i=f(x_i)-f(x_{i-1})$. Maka
\[
U(\mathcal{P},f)-L(\mathcal{P},f)
=\Delta x\sum_{i=1}^{n}\bigl(f(x_i)-f(x_{i-1})\bigr)
=\frac{(b-a)\bigl(f(b)-f(a)\bigr)}{n},
\]
karena jumlah di tengah teleskopik. Pilih $n>(b-a)(f(b)-f(a))/\varepsilon$ dan kriteria
terpenuhi. Bila $f$ turun, terapkan hasil ini pada $-f$ dan pakai linearitas.
\end{proof}

Bukti itu tidak memakai $\delta$: yang dipakai hanya bahwa jumlah semua lompatan tidak
melebihi $f(b)-f(a)$, sehingga fungsi monoton mudah diintegralkan meski grafiknya penuh
sudut.

Kedua kelas itu dapat digabung: fungsi yang kontinu kecuali di berhingga banyak titik
tetap terintegralkan (Rudin, Teorema 6.10, tidak diulang di sini).

Sekarang sisi lain: fungsi terbatas yang gagal. Contoh paling bersih adalah fungsi
Dirichlet, bernilai $1$ pada bilangan rasional dan $0$ pada bilangan irasional.

\begin{contoh}[Fungsi Dirichlet tidak terintegralkan]\index{fungsi dirichlet}
Misalkan $D(x)=1$ bila $x$ rasional dan $D(x)=0$ bila $x$ irasional, pada $[0,1]$.
Setiap potongan memuat keduanya (Bab 3, kerapatan $\mathbb{Q}$), sehingga $M_i=1$ dan
$m_i=0$ untuk setiap $i$, dan
\[
U(\mathcal{P},D)=\sum_{i=1}^{n}1\cdot\Delta x_i=1,
\qquad
L(\mathcal{P},D)=0
\]
untuk \emph{setiap} partisi. Selisihnya selalu $1$, tidak pernah lebih kecil dari
$\varepsilon=1/2$, jadi menurut kriteria $D$ tidak terintegralkan meskipun terbatas.
\end{contoh}

Kegagalan itu bukan karena $D$ terputus-putus di banyak titik, melainkan karena
lompatannya tersebar di mana-mana. Fungsi berikut menunjukkan bahwa banyaknya titik
diskontinu bukan penghalang.

\begin{teorema}[Fungsi Thomae]\index{fungsi thomae}
Definisikan $T$ pada $[0,1]$ dengan $T(x)=1/q$ bila $x=p/q$ dalam bentuk paling
sederhana dengan $q\ge1$, dan $T(x)=0$ bila $x$ irasional. Maka $T$ terintegralkan
Riemann dan $\int_0^1 T=0$.
\end{teorema}

\begin{proof}
Ambil $\varepsilon>0$. Himpunan $S=\{x:T(x)\ge\varepsilon\}$ berisi pecahan dengan
penyebut $q\le1/\varepsilon$; untuk setiap $q$ hanya ada $q+1$ pecahan seperti itu di
dalam $[0,1]$, sehingga $S$ berhingga. Tutup $S$ dengan berhingga banyak selang
terbuka yang panjang totalnya kurang dari $\varepsilon$.

Bentuk partisi yang memuat semua ujung selang penutup itu. Potongan yang memuat $S$
panjangnya berjumlah kurang dari $\varepsilon$ dan di sana $M_i-m_i\le1$, jadi
sumbangannya kurang dari $\varepsilon$. Potongan lain memenuhi $M_i\le\varepsilon$ dan
$m_i\ge0$, sumbangannya paling banyak $\varepsilon$. Karena itu
$U(\mathcal{P},T)-L(\mathcal{P},T)<2\varepsilon$, dan kriteria menutup bukti. Nilainya
nol sebab setiap potongan memuat bilangan irasional, sehingga $L(\mathcal{P},T)=0$.
\end{proof}

Jadi $T$ terputus di setiap bilangan rasional, yang rapat di $[0,1]$, tetapi $T$ tetap
terintegralkan. Yang menentukan bukan jumlah titik diskontinu, melainkan besar ruang
yang ditempatinya: kriteria Lebesgue menyatakan fungsi terbatas terintegralkan Riemann
jika dan hanya jika titik diskontinunya berukuran nol (Rudin, Teorema 11.33(b); Abbott,
Bab 7). Kita mengutipnya tanpa bukti karena ukuran menuntut teori ukur.

\begin{tantangan}
Definisikan $h$ pada $[0,1]$ dengan $h(x)=1$ bila $x=p/q$ dalam bentuk paling sederhana
dan penyebut $q$ ganjil, serta $h(x)=0$ untuk semua $x$ lain, termasuk setiap bilangan
irasional. Apakah $h$ terintegralkan Riemann pada $[0,1]$? Bila ya, berapa nilai
integralnya? Bila tidak, buktikan memakai kriteria keterintegralan, bukan dengan
menghitung titik.

Petunjuk 1: jangan menghitung banyaknya titik tempat $h$ bernilai satu. Banyaknya tak
berhingga, dan kriteria tidak meminta hitungan itu.

Petunjuk 2: ambil sebarang potongan $[x_{i-1},x_i]$ dengan $x_{i-1}<x_i$. Potongan itu
memuat pecahan berpenyebut ganjil, jadi $M_i=1$, dan memuat bilangan irasional, jadi
$m_i=0$. Dengan begitu $U-L=1$ untuk setiap partisi.

Kalau mentok, kembali ke bagian Kelas fungsi yang terintegralkan, contoh Dirichlet.
\end{tantangan}

\kamus{keterintegralan}{Riemann integrability}{Sifat fungsi yang taksiran bawah dan
atasnya dapat dirapatkan sampai selisihnya hilang.}

\begin{untukanak}
Pernahkah kita melihat speedometer mobil, lalu bertanya berapa jarak yang sudah
ditempuh? Jarak tidak ditampilkan di situ, hanya kecepatan sesaat. Tetapi kalau kita
mencatat kecepatan setiap detik, kita bisa menjumlahkannya: satu detik dengan
kecepatan 10 meter per detik berarti 10 meter, detik berikutnya 12 meter, dan
seterusnya. Jumlah semua potongan itu adalah jarak.

Masalahnya, kecepatan berubah di tengah detik. Mobil bisa mempercepat lalu melambat
dalam satu detik yang sama. Kalau kita memakai kecepatan di awal detik, jarak yang
kita hitung terlalu kecil. Kalau memakai kecepatan di akhir detik, hasilnya terlalu
besar. Yang benar ada di antaranya.

Kalau kita mencatat sepuluh kali per detik, kedua taksiran makin dekat; seratus kali,
lebih dekat lagi. Matematikawan bertanya lebih jauh: kalau pencatatan bisa serapat apa
pun, apakah kedua taksiran itu menuju satu angka yang sama? Jawabannya tidak selalu ya,
dan di bab ini kita akan tahu kapan ya dan mengapa.
\end{untukanak}

\section{Integral sebagai fungsi akumulasi}
\tingkat{menengah}

Sampai sini integral hanyalah bilangan yang melekat pada selang. Kita dapat
memandangnya sebagai fungsi dengan membiarkan ujung atas bergerak: integralnya menjadi
laporan tentang berapa banyak yang sudah terkumpul sampai titik itu.

\begin{definisi}[Fungsi akumulasi]\index{fungsi akumulasi}
Misalkan $f$ terintegralkan pada $[a,b]$. Untuk $x\in[a,b]$ didefinisikan
\[
F(x)=\int_a^x f(t)\,dt .
\]
Untuk $x<a$ dipakai kesepakatan $\int_a^x f=-\int_x^a f$, dan untuk $x=a$ nilainya
nol. Fungsi $F$ disebut \emph{fungsi akumulasi} dari $f$ pada $[a,b]$.
\end{definisi}

Dengan kesepakatan itu sifat aditif berlaku untuk urutan titik apa pun, sehingga
\begin{equation}
F(y)-F(x)=\int_x^{y} f(t)\,dt
\end{equation}
untuk semua $x,y\in[a,b]$. Inilah bentuk yang dipakai pada semua bukti berikutnya:
pertambahan tabungan antara dua waktu adalah integral laju pemasukan di antaranya.

Fungsi akumulasi selalu lebih baik daripada fungsi asalnya: ia kontinu walaupun $f$
tidak, dan terdiferensialkan di setiap titik tempat $f$ kontinu.

\begin{proposisi}[Akumulasi selalu kontinu]
Jika $f$ terbatas pada $[a,b]$ dengan $|f|\le M$ dan terintegralkan, maka $F$ yang
didefinisikan di atas memenuhi
\[
|F(y)-F(x)|\le M\,|y-x|
\]
untuk semua $x,y\in[a,b]$. Khususnya $F$ kontinu pada $[a,b]$.
\end{proposisi}

\begin{proof}
Untuk $x<y$, sifat monoton dan bagian (d) Proposisi sifat dasar memberi
\[
|F(y)-F(x)|=\Bigl|\int_x^{y}f\Bigr|\le\int_x^{y}|f|\le M(y-x).
\]
Untuk $y<x$ hasil yang sama dengan tanda berlawanan. Jadi $F$ memenuhi syarat
Lipschitz dengan konstanta $M$, dan setiap fungsi Lipschitz kontinu.
\end{proof}

Sekarang arah pertama hubungan turunan-integral. Kalimat sehari-harinya: laju
pertumbuhan tabungan pada suatu saat sama dengan laju pemasukan pada saat itu.

\begin{teorema}[Teorema dasar kalkulus, bagian pertama]\index{teorema dasar kalkulus}
Misalkan $f$ terintegralkan pada $[a,b]$ dan kontinu di titik $x_0\in(a,b)$. Maka
fungsi akumulasi $F$ terdiferensialkan di $x_0$ dan
\[
F'(x_0)=f(x_0).
\]
\end{teorema}

\begin{proof}
Ambil $h\ne0$ dengan $x_0+h\in[a,b]$. Dari bentuk pertambahan,
\[
\frac{F(x_0+h)-F(x_0)}{h}-f(x_0)
=\frac{1}{h}\int_{x_0}^{x_0+h}\bigl(f(t)-f(x_0)\bigr)\,dt ,
\]
sebab integral fungsi konstan $f(x_0)$ pada selang panjang $|h|$ bernilai
$f(x_0)h$.

Ambil $\varepsilon>0$. Karena $f$ kontinu di $x_0$, ada $\delta>0$ sehingga
$|t-x_0|<\delta$ berakibat $|f(t)-f(x_0)|<\varepsilon$. Untuk $0<|h|<\delta$ seluruh
selang integrasi berada di dalam $(x_0-\delta,x_0+\delta)$, sehingga dengan sifat
monoton dan bagian (d) Proposisi sifat dasar,
\[
\left|\frac{1}{h}\int_{x_0}^{x_0+h}\bigl(f(t)-f(x_0)\bigr)dt\right|
\le\frac{1}{|h|}\left|\int_{x_0}^{x_0+h}\bigl|f(t)-f(x_0)\bigr|\,dt\right|
\le\frac{1}{|h|}\,\varepsilon\,|h|=\varepsilon .
\]
Karena $\varepsilon$ sebarang, hasil bagi selisih itu menuju $f(x_0)$.
\end{proof}

Teorema itu hanya menuntut kekontinuan di satu titik. Di tempat $f$ melompat, turunan
$F$ boleh gagal ada, dan itu bukan kelemahan bukti: untuk $f(t)=\operatorname{sgn} t$
pada $[-1,1]$ berlaku $F(x)=\int_0^x\operatorname{sgn} t\,dt=|x|$, yang tidak
terdiferensialkan di $0$, persis di tempat tanda $f$ berubah.

Satu akibat penting: setiap fungsi kontinu punya antiturunan, dan antiturunannya dapat
ditulis sebagai integral. Bab 12 menurunkan aturan turunan tetapi tidak menjawab apakah
setiap fungsi kontinu merupakan turunan dari sesuatu. Sekarang pertanyaan itu terjawab,
dan jawabannya konstruktif.

\begin{akibat}[Antiturunan selalu ada]
Jika $f$ kontinu pada $[a,b]$, maka $F(x)=\int_a^x f$ terdiferensialkan pada $[a,b]$
dengan $F'=f$. Jadi $F$ adalah antiturunan $f$ pada $[a,b]$.
\end{akibat}

\begin{proof}
Kekontinuan $f$ memberi keterintegralan (Akibat sebelumnya) dan kekontinuan di setiap
titik, sehingga teorema bagian pertama berlaku di setiap $x_0\in(a,b)$. Di titik ujung
dipakai turunan satu sisi.
\end{proof}

Aturan rantai memberi satu keluwesan tambahan: bila batas atas bergantung pada $x$,
\[
\frac{d}{dx}\int_a^{g(x)}f(t)\,dt=f\bigl(g(x)\bigr)\,g'(x),
\]
asalkan $g$ terdiferensialkan. Contohnya,
$\frac{d}{dx}\int_0^{x^2}e^{-t^2}dt=2x\,e^{-x^4}$.

Kedua arah itu mudah tertukar, jadi ada baiknya digambar sekali.
Gambar~\ref{fig:tdk-dua-arah} menaruh keduanya pada satu diagram: panah atas
mengumpulkan potongan tipis, panah bawah menurunkan. Yang dijamin teorema adalah bahwa
menurunkan hasil pengumpulan mengembalikan fungsi semula, bukan kebalikan aljabar yang
rapi, sebab antiturunan hanya pasti sampai konstanta.

\begin{figure}[htbp]
\centering
\begin{tikzpicture}[>=Stealth,every node/.style={font=\footnotesize},
  kotak/.style={draw,rounded corners=3pt,align=center,inner sep=5pt,
    minimum height=10mm,text width=3.3cm}]
\node[kotak,fill=layerindigo!8,draw=layerindigo] (f) at (0,0) {fungsi $f$};
\node[kotak,fill=layerbrown!8,draw=layerbrown] (F) at (7.6,0)
  {fungsi akumulasi $F$};
\draw[->,line width=0.9pt,draw=layerindigo] ([yshift=2mm]f.north) to[bend left=24]
  node[above,text=layerindigo,align=center]
  {integral tentu: kumpulkan potongan tipis} ([yshift=2mm]F.north);
\draw[->,line width=0.9pt,draw=layerbrown] ([yshift=-2mm]F.south) to[bend left=24]
  node[below,text=layerbrown,align=center]
  {turunan: $F'=f$, inilah bagian 1} ([yshift=-2mm]f.south);
\node[draw=layergrey,dashed,rounded corners=3pt,align=center,inner sep=4pt]
  at (3.8,-2.7)
  {bagian 2 menutup lingkaran: $\int_a^b F'(x)\,dx=F(b)-F(a)$};
\end{tikzpicture}
\caption{Dua arah teorema dasar kalkulus sebagai panah bolak-balik. Panah atas:
mengumpulkan potongan tipis menghasilkan $F(x)=\int_a^x f$. Panah bawah: menurunkan $F$
mengembalikan $f$, dan itulah bagian pertama. Kotak putus-putus di bawah adalah bagian
kedua, yang memakai antiturunan untuk menghitung integral. Untuk $f(x)=x^2$ pada
$[0,1]$, $F(x)=x^3/3$, dan bagian kedua memberi $F(1)-F(0)=1/3$, angka yang sama dengan
pengurungan pada Gambar~\ref{fig:rect-halus}.}
\label{fig:tdk-dua-arah}
\end{figure}

\section{Teorema dasar kalkulus, bagian kedua}
\tingkat{lanjut}

Bagian pertama berjalan dari integral ke turunan. Bagian kedua berjalan sebaliknya:
bila kita sudah memegang antiturunan, integral dapat dihitung dengan satu pengurangan.
Kalimat sehari-harinya: total perubahan sama dengan nilai akhir dikurangi nilai awal.

\begin{teorema}[Teorema dasar kalkulus, bagian kedua]
Misalkan $F$ kontinu pada $[a,b]$, terdiferensialkan pada $(a,b)$, dan $F'$
terintegralkan Riemann pada $[a,b]$. Maka
\[
\int_a^b F'(x)\,dx=F(b)-F(a).
\]
\end{teorema}

\begin{proof}
Ambil partisi $\mathcal{P}=\{x_0,\dots,x_n\}$ dari $[a,b]$. Pada setiap subselang, $F$
kontinu pada $[x_{i-1},x_i]$ dan terdiferensialkan di dalamnya, sehingga teorema nilai
rata-rata (Bab 12) memberi $t_i\in(x_{i-1},x_i)$ dengan
\[
F(x_i)-F(x_{i-1})=F'(t_i)\,\Delta x_i .
\]
Karena $m_i\le F'(t_i)\le M_i$, penjumlahan atas $i$ memberi
\[
L(\mathcal{P},F')\le\sum_{i=1}^{n}\bigl(F(x_i)-F(x_{i-1})\bigr)\le U(\mathcal{P},F').
\]
Jumlah di tengah teleskopik dan sama dengan $F(b)-F(a)$, jadi bilangan itu terkurung di
antara jumlah bawah dan jumlah atas setiap partisi. Karena $F'$ terintegralkan, kedua
kurungan menuju nilai yang sama, yaitu $\int_a^b F'$.
\end{proof}

\begin{resep}
\begin{enumerate}
\item Cari $F$ dengan $F'=f$, lalu periksa dengan menurunkan $F$. Bila turunannya tidak
      kembali ke $f$, calonnya salah dan langkah berikutnya sia-sia.
\item Tulis $F$ pada batas atas, tulis $F$ pada batas bawah, lalu kurangkan:
      $F(b)-F(a)$. Jangan sisipkan $+C$; konstanta itu saling menghapus.
\item Rapikan tandanya bila selangnya terbalik: $\int_b^a f=-\int_a^b f$, dan
      $\int_a^a f=0$.
\item Uji dengan satu potongan kasar, trapesium atau \texttt{quad}. Bila hasilnya jauh,
      galatnya ada di langkah 1 atau 2, bukan di langkah 3.
\end{enumerate}
Peringatan urutan: memeriksa antiturunan harus terjadi sebelum mengurangkan. Dua nilai
fungsi sembarang tetap memberi selisih yang rapi dan masuk akal, sehingga keliru di
langkah 1 tidak terlihat dari bentuk jawabannya.
\end{resep}

Begitu pula kedua bagian itu tidak saling menggantikan: bagian pertama menurunkan
fungsi akumulasi tanpa memerlukan antiturunan, sedangkan bagian kedua memerlukan
antiturunan dan keterintegralan turunannya.

Bagian kedua menjelaskan mengapa hasilnya disebut perubahan bersih, bukan luas. Untuk
$F(x)=\frac{x^4}{4}-\frac{x^2}{2}$ pada $[-1,1]$, $F'(x)=x^3-x$, sehingga
\[
\int_{-1}^{1}(x^3-x)\,dx=F(1)-F(-1)=\Bigl(\frac14-\frac12\Bigr)-\Bigl(\frac14-\frac12\Bigr)=0 .
\]
Daerah di atas sumbu seluas $1/4$, daerah di bawah sumbu juga $1/4$, jadi luas totalnya
$1/2$ sedangkan integralnya nol karena tandanya berlawanan.
Gambar~\ref{fig:luas-bertanda} memperlihatkan keduanya.

\begin{figure}[htbp]
\centering
\begin{tikzpicture}
\begin{axis}[width=0.74\linewidth,height=5.4cm,
  axis lines=middle,xmin=-1.3,xmax=1.3,ymin=-0.6,ymax=0.6,
  xtick={-1,0,1},ytick=\empty,
  tick label style={font=\scriptsize},
  xlabel={$x$},ylabel={$y=x^3-x$},
  xlabel style={font=\scriptsize},ylabel style={font=\scriptsize,rotate=-90}]
\addplot[fill=layerjade!22,draw=layerjade!70,domain=-1:0,samples=60]
  {x^3-x} \closedcycle;
\addplot[fill=layerplum!22,draw=layerplum!70,domain=0:1,samples=60]
  {x^3-x} \closedcycle;
\addplot[thick,ink,domain=-1.2:1.2,samples=120] {x^3-x};
\node[font=\scriptsize,text=layerjade!55!black] at (axis cs:-0.52,0.36)
  {$+\frac14$};
\node[font=\scriptsize,text=layerplum!55!black] at (axis cs:0.52,-0.36)
  {$-\frac14$};
\end{axis}
\end{tikzpicture}
\caption{Luas bertanda untuk $F'(x)=x^3-x$ pada $[-1,1]$. Daerah di atas sumbu-$x$
menyumbang $+1/4$, daerah di bawah sumbu menyumbang $-1/4$, sehingga
$\int_{-1}^{1}(x^3-x)\,dx=F(1)-F(-1)=0$. Luas geometrinya tetap $1/2$; yang nol adalah
jumlah bertandanya, dan itulah yang dihitung teorema bagian kedua.}
\label{fig:luas-bertanda}
\end{figure}

Hipotesis keterintegralan pada bagian kedua tidak berlebihan. Ada fungsi
terdiferensialkan yang turunannya terbatas tetapi tidak Riemann-terintegralkan; contoh
klasiknya ditemukan Volterra pada abad kesembilan belas, dan kita tidak
membuktikannya di sini. Teorema bagian kedua bukan pernyataan tanpa syarat.

\section{Substitusi dan integral parsial}
\tingkat{lanjut}

Dua aturan hitung integral yang paling sering dipakai bukan aksioma baru: keduanya
aturan turunan yang dibalik, dan pembuktiannya hanya memakai teorema dasar kalkulus.
Yang pertama membalik aturan rantai.

\begin{proposisi}[Substitusi]\index{substitusi}
Misalkan $\varphi$ terdiferensialkan dengan $\varphi'$ kontinu pada $[a,b]$, dan $f$
kontinu pada selang yang memuat $\varphi([a,b])$. Maka
\[
\int_a^b f\bigl(\varphi(t)\bigr)\,\varphi'(t)\,dt
=\int_{\varphi(a)}^{\varphi(b)}f(u)\,du .
\]
\end{proposisi}

\begin{proof}
Karena $f$ kontinu, ia mempunyai antiturunan $G(u)=\int_c^u f$, dengan $c$ di dalam
selang nilai $\varphi$; menurut teorema bagian pertama, $G'=f$. Aturan rantai (Bab 12)
memberi
\[
\frac{d}{dt}G\bigl(\varphi(t)\bigr)=G'\bigl(\varphi(t)\bigr)\varphi'(t)
=f\bigl(\varphi(t)\bigr)\varphi'(t).
\]
Ruas kanan kontinu, jadi terintegralkan. Teorema bagian kedua pada fungsi
$G\circ\varphi$ memberi
\[
\int_a^b f\bigl(\varphi(t)\bigr)\varphi'(t)\,dt
=G\bigl(\varphi(b)\bigr)-G\bigl(\varphi(a)\bigr)
=\int_{\varphi(a)}^{\varphi(b)}f(u)\,du .
\]
\end{proof}

Aturan kedua membalik aturan perkalian.

\begin{proposisi}[Integral parsial]\index{integral parsial}
Misalkan $F$ dan $G$ terdiferensialkan pada $[a,b]$ dengan $F'$ dan $G'$ kontinu. Maka
\[
\int_a^b F(x)\,G'(x)\,dx=F(b)G(b)-F(a)G(a)-\int_a^b F'(x)\,G(x)\,dx .
\]
\end{proposisi}

\begin{proof}
Misalkan $H=FG$. Aturan perkalian memberi $H'=F'G+FG'$, dan ruas kanan kontinu
sehingga $H'$ terintegralkan. Teorema bagian kedua pada $H$ memberi
\[
F(b)G(b)-F(a)G(a)=\int_a^b H'=\int_a^b F'G+\int_a^b FG' .
\]
Pindahkan suku pertama ke ruas kiri.
\end{proof}

\kamus{antiturunan}{antiderivative}{Fungsi yang turunannya adalah fungsi yang sedang
kita integralkan, seperti arah pulang dari sebuah jalan: ia hanya berguna kalau arah
berangkatnya sudah diketahui.}

Sebagai pemeriksaan angka, ambil $\int_0^1 x\,e^{x^2}dx$. Substitusi $u=x^2$,
$du=2x\,dx$ memberi
\[
\int_0^1 x\,e^{x^2}dx=\frac12\int_0^1 e^{u}du=\frac{e-1}{2}=0.8591409142295225 .
\]
\texttt{scipy.integrate.quad} memberi $0.8591409142295227$, selisih $2\times10^{-16}$
dari pembulatan mesin. Untuk aturan kedua, $F(x)=x$ dan $G'(x)=e^x$ memberi
$\int_0^1 xe^xdx=e-(e-1)=1$, dan \texttt{quad} mengembalikan $1.0$ dengan perkiraan
galat $1.1\times10^{-14}$.

\section{Teorema nilai rata-rata untuk integral}
\tingkat{menengah}

Rata-rata biasanya dihitung dengan menjumlah lalu membagi banyaknya data. Untuk fungsi
yang berubah terus, penjumlahan diganti integral dan pembaginya panjang selang: rata-rata
adalah tinggi persegi panjang yang luasnya sama dengan luas di bawah kurva.

\begin{definisi}[Nilai rata-rata]
Untuk $f$ terintegralkan pada $[a,b]$, \emph{nilai rata-rata} $f$ pada selang itu
adalah
\[
\bar{f}=\frac{1}{b-a}\int_a^b f(x)\,dx .
\]
\end{definisi}

\begin{teorema}[Nilai rata-rata untuk integral]\index{nilai rata-rata integral}
Jika $f$ kontinu pada $[a,b]$, maka terdapat $c\in[a,b]$ dengan
\[
f(c)=\frac{1}{b-a}\int_a^b f(x)\,dx .
\]
\end{teorema}

\begin{proof}
Karena $f$ kontinu pada selang tertutup terbatas, ia mencapai nilai terkecil $m$ dan
nilai terbesar $M$ di sana (Bab 11). Sifat monoton memberi
\[
m(b-a)\le\int_a^b f\le M(b-a),
\]
sehingga $\bar{f}\in[m,M]$. Menurut teorema nilai antara (Bab 11), fungsi kontinu $f$
mengambil setiap nilai di antara $m$ dan $M$, jadi ada $c\in[a,b]$ dengan
$f(c)=\bar{f}$.
\end{proof}

Sebagai angka, untuk $f(x)=x^2$ pada $[0,1]$ rata-ratanya $1/3$, dan titik yang
dicapai teorema adalah $c=1/\sqrt3=0.5773502691896258$, sebab
$c^2=0.3333333333333334$. Titik itu bukan titik tengah selang: teorema menjamin
keberadaan, bukan lokasi.

\section{Contoh hitung}
\tingkat{menengah}

Sekarang kita kerjakan tiga hal: jumlah atas dan bawah $f(x)=x^2$ pada $[0,1]$,
perbandingan lima aturan kuadratur, dan satu integral tanpa antiturunan elementer.
Pecahan dihitung dengan \texttt{fractions.Fraction}, tanpa pembulatan.

Mulai dari $n$ potong sama lebar dengan $x_i=i/n$ dan $\Delta x=1/n$, sehingga
\[
U_n=\sum_{i=1}^{n}\Bigl(\frac{i}{n}\Bigr)^2\frac1n=\frac{1}{n^3}\sum_{i=1}^{n}i^2
=\frac{(n+1)(2n+1)}{6n^2}=\frac13+\frac{1}{2n}+\frac{1}{6n^2},
\]
memakai rumus jumlah kuadrat dari Bab 1. Dengan cara yang sama
$L_n=\frac13-\frac{1}{2n}+\frac{1}{6n^2}$, sehingga
\begin{equation}
U_n-L_n=\frac1n .
\end{equation}
Hasil itu bukan hampiran. Untuk setiap $n$ selisihnya persis $1/n$; karena $1/n$ dapat
dibuat lebih kecil dari $\varepsilon$ mana pun dengan $n>1/\varepsilon$, kriteria
keterintegralan terpenuhi. Jadi tabel berikut sekaligus bukti bahwa
$\int_0^1x^2dx=1/3$: kedua jumlah mengapit $1/3$, dan jaraknya menuju nol.

\begin{center}\small
\begin{tabular}{rrrr}
\hline
$n$ & $U_n$ & $L_n$ & $U_n-L_n$ \\
\hline
1 & $1.0000000000$ & $0.0000000000$ & $1$ \\
4 & $0.4687500000$ & $0.2187500000$ & $1/4$ \\
16 & $0.3652343750$ & $0.3027343750$ & $1/16$ \\
64 & $0.3411865234$ & $0.3255615234$ & $1/64$ \\
256 & $0.3352890015$ & $0.3313827515$ & $1/256$ \\
1024 & $0.3338217735$ & $0.3328452110$ & $1/1024$ \\
\hline
\end{tabular}
\end{center}

Pada $n=1024$ selisihnya $9.77\times10^{-4}$ dan jarak $U_n$ ke $1/3$ separuh dari itu.
Setiap kali $n$ dilipatgandakan selisihnya tepat separuh, dan itulah sebabnya jumlah
atas dan bawah bukan cara efisien untuk mencapai ketelitian tinggi.

\begin{tangan}
Bahan: satu lembar kertas berpetak, gunting, penggaris, selotip, pensil.

Langkah. (1) Di kertas berpetak gambar kurva yang naik makin cepat: satu petak dari
tepi kiri tingginya satu petak, dua petak tingginya empat petak, tiga petak tingginya
sembilan petak. (2) Gunting daerah di bawah kurva menjadi jalur tegak selebar satu
petak. (3) Susun jalur itu berdampingan di meja, dari terpendek ke terpanjang, alas
rata. (4) Perhatikan sisi atasnya: yang muncul tangga bergerigi, bukan garis mulus.
(5) Ulangi dengan jalur selebar setengah petak, lalu seperempat petak.

Yang harus kita lihat: gerigi di sisi atas makin halus ketika jalur makin tipis,
sedangkan luas totalnya tidak berubah, sebab memotong lalu menyusun ulang tidak
menambah atau mengurangi kertas. Jadi luas daerah melengkung itu sama dengan luas
bangun bergerigi, dan bangun bergerigi itu makin mirip satu persegi panjang. Tinggi
persegi panjang itulah yang nanti kita sebut nilai rata-rata fungsi pada selang itu.
\end{tangan}

Sekarang bandingkan lima aturan kuadratur pada integral yang sama: ujung kiri, ujung
kanan, titik tengah $x_i=(i-\frac12)/n$, trapesium, dan Simpson dengan bobot
$1,4,2,4,\dots,4,1$ dibagi $3n$. Galat dihitung terhadap nilai eksak $1/3$.

\begin{center}\small
\begin{tabular}{rrrrrr}
\hline
$n$ & kiri & kanan & titik tengah & trapesium & Simpson \\
\hline
10 & $4.833\times10^{-2}$ & $5.167\times10^{-2}$ & $8.333\times10^{-4}$
   & $1.667\times10^{-3}$ & $1.1\times10^{-16}$ \\
100 & $4.983\times10^{-3}$ & $5.017\times10^{-3}$ & $8.333\times10^{-6}$
   & $1.667\times10^{-5}$ & $0$ \\
1000 & $4.998\times10^{-4}$ & $5.002\times10^{-4}$ & $8.333\times10^{-8}$
   & $1.667\times10^{-7}$ & $0$ \\
\hline
\end{tabular}
\end{center}

Tiga pola terbaca. Galat aturan kiri dan kanan turun sepuluh kali setiap $n$ naik
sepuluh kali: orde $\mathcal{O}(1/n)$, sekitar $1/(2n)$ karena $f$ naik. Galat titik
tengah dan trapesium turun seratus kali, orde $\mathcal{O}(1/n^2)$, dengan nilai
$1/(12n^2)$ dan $1/(6n^2)$; keduanya berlawanan tanda dan trapesium dua kali lipat,
sehingga $(2M+T)/3$ menghapus sukunya. Simpson tepat sampai pembulatan mesin, sebab
aturannya eksak untuk polinomial berderajat paling tinggi tiga.

Integral ketiga tidak punya antiturunan elementer:
\[
\int_0^1 e^{-x^2}dx .
\]
Ketiadaan itu akibat teorema Liouville, kita kutip tanpa bukti. Acuannya fungsi galat
$\frac{\sqrt{\pi}}{2}\operatorname{erf}(1)=0.746824132812427$; \texttt{quad} memberi
angka yang sama dengan perkiraan galat $8.3\times10^{-15}$. Aturan Simpson memberi galat
$3.563\times10^{-4}$ pada $n=2$, $1.988\times10^{-6}$ pada $n=8$, dan
$7.795\times10^{-9}$ pada $n=32$.

\kertas{Untuk $f(x)=x^2$ pada $[0,1]$, turunkan sendiri rumus
$U_n-L_n=1/n$ dari rumus jumlah kuadrat, lalu hitung $U_{10}$ dan $L_{10}$ sebagai
pecahan. Periksa bahwa rata-ratanya $(U_{10}+L_{10})/2$ sama dengan
$1/3+1/(6\cdot100)$, bukan $1/3$. Jelaskan mengapa rata-rata itu tetap bukan taksiran
terbaik, padahal galatnya jauh lebih kecil daripada galat $U_{10}$ sendiri.}

\section{Integral tak wajar dan kaitannya dengan deret}
\tingkat{lanjut}

Selang integrasi sejauh ini selalu terbatas dan fungsinya terbatas. Banyak perhitungan
nyata tidak begitu: kita ingin luas daerah yang memanjang tanpa ujung, atau integral
fungsi yang meledak di ujung selang. Caranya bukan memperluas definisi Riemann,
melainkan menulis integral itu sebagai limit integral biasa.

\begin{definisi}[Integral tak wajar]\index{integral tak wajar}
Bila $f$ terintegralkan Riemann pada $[a,R]$ untuk setiap $R>a$, didefinisikan
\[
\int_a^\infty f=\lim_{R\to\infty}\int_a^R f ,
\]
asalkan limit itu ada dan berhingga. Bila $f$ terintegralkan pada $[c,b]$ untuk setiap
$c\in(a,b]$ tetapi tidak terbatas di dekat $a$, didefinisikan
\[
\int_a^b f=\lim_{c\to a^{+}}\int_c^b f .
\]
Dalam kedua kasus, integral disebut \emph{konvergen} bila limitnya ada, dan
\emph{berdivergen} bila tidak.
\end{definisi}

Contoh terpenting adalah keluarga $f(x)=x^{-p}$ pada $[1,\infty)$. Untuk $p\ne1$,
\[
\int_1^R x^{-p}dx=\frac{R^{1-p}-1}{1-p},
\]
sehingga limitnya ada hanya bila $p>1$, dengan nilai $1/(p-1)$; untuk $p=1$ hasilnya
$\ln R$, yang menuju tak berhingga. Jadi $\int_1^\infty x^{-2}dx=1$, dan integral
parsialnya menuju nilai itu dengan laju $1/R$: $0.9$, $0.99$, $0.999$ untuk
$R=10,100,1000$.

Kaitan dengan deret Bab 5 bukan kebetulan: untuk $f$ positif dan turun, luas persegi
panjang setinggi $f(n)$ mengurung luas di bawah kurva, sehingga
\[
\sum_{n=N+1}^{\infty}f(n)\ \le\ \int_N^\infty f(x)\,dx\ \le\ \sum_{n=N}^{\infty}f(n) .
\]
Karena itu uji integral berlaku: $\sum f(n)$ konvergen jika dan hanya jika
$\int_1^\infty f$ konvergen. Untuk $f(x)=1/x^2$ keduanya konvergen, dengan
$\sum_{n=1}^{\infty}1/n^2=\pi^2/6=1.6449340668482264$. Ketaksamaan kiri juga memberi
batas galat, dan makin besar $N$ makin dekat galat ke batas itu:

\begin{center}\small
\begin{tabular}{rrrr}
\hline
$N$ & galat jumlah parsial & $1/N$ \\
\hline
10 & $9.516634\times10^{-2}$ & $0.1$ \\
100 & $9.950167\times10^{-3}$ & $0.01$ \\
1000 & $9.995002\times10^{-4}$ & $0.001$ \\
\hline
\end{tabular}
\end{center}

Perlu ditegaskan: integral tak wajar bukan integral Riemann, melainkan limitnya. Karena
itu substitusi dan integral parsial hanya berlaku bila limitnya ada, dan untuk fungsi
berganti tanda kekonvergenannya belum tentu mutlak.

\section{Lab komputasi}
\tingkat{lanjut}

Kode berikut dijalankan dengan Python 3, \texttt{scipy} 1.18.0, dan modul
\texttt{fractions} untuk aritmetika eksak. Semua partisi sama lebar pada $[0,1]$.

\begin{lab}
\begin{verbatim}
import math
from fractions import Fraction as F
from scipy.integrate import quad

def jumlah(n):
    U = sum(F(i * i, n ** 3) for i in range(1, n + 1))
    L = sum(F(i * i, n ** 3) for i in range(n))
    return U, L

print("n        U_n          L_n          U_n - L_n")
for n in (1, 2, 4, 8, 16, 32, 64, 128, 256, 512, 1024):
    U, L = jumlah(n)
    print(f"{n:5d}  {float(U):.10f}  {float(L):.10f}  {float(U - L):.2e}")

def aturan(f, n):
    h = 1 / n
    kiri = sum(f(i * h) for i in range(n)) * h
    kanan = sum(f(i * h) for i in range(1, n + 1)) * h
    tengah = sum(f((i + 0.5) * h) for i in range(n)) * h
    trap = (f(0) / 2 + sum(f(i * h) for i in range(1, n)) + f(1) / 2) * h
    simp = (f(0) + f(1) + sum((4 if i % 2 else 2) * f(i * h)
                              for i in range(1, n))) * h / 3
    return kiri, kanan, tengah, trap, simp

f = lambda x: x * x
print()
print("galat |hampiran - 1/3|, integral x^2 pada [0,1]")
print("   n      kiri      kanan     tengah   trapesium    Simpson")
for n in (10, 100, 1000):
    galat = [f"{abs(v - 1 / 3):.3e}" for v in aturan(f, n)]
    print(f"{n:4d}  " + "  ".join(galat))

g = lambda x: math.exp(-x * x)
v, e = quad(g, 0, 1)
acuan = math.sqrt(math.pi) / 2 * math.erf(1)
print()
print(f"quad exp(-x^2) = {v:.15f}   galat perkiraan {e:.1e}")
print(f"acuan erf      = {acuan:.15f}   selisih {v - acuan:+.1e}")
print(f"quad x^2       = {quad(f, 0, 1)[0]:.17f}")
print()
print("galat |Simpson - acuan| untuk exp(-x^2)")
for n in (2, 4, 8, 16, 32):
    print(f"  n={n:3d}  {abs(aturan(g, n)[4] - acuan):.3e}")
print()
print("galat |trapesium - acuan| untuk exp(-x^2)")
for n in (10, 100, 1000):
    print(f"  n={n:4d}  {abs(aturan(g, n)[3] - acuan):.3e}")
\end{verbatim}
\end{lab}

Keluaran yang diperoleh, apa adanya:

\begin{verbatim}
n        U_n          L_n          U_n - L_n
    1  1.0000000000  0.0000000000  1.00e+00
    2  0.6250000000  0.1250000000  5.00e-01
    4  0.4687500000  0.2187500000  2.50e-01
    8  0.3984375000  0.2734375000  1.25e-01
   16  0.3652343750  0.3027343750  6.25e-02
   32  0.3491210938  0.3178710938  3.12e-02
   64  0.3411865234  0.3255615234  1.56e-02
  128  0.3372497559  0.3294372559  7.81e-03
  256  0.3352890015  0.3313827515  3.91e-03
  512  0.3343105316  0.3323574066  1.95e-03
 1024  0.3338217735  0.3328452110  9.77e-04

galat |hampiran - 1/3|, integral x^2 pada [0,1]
   n      kiri      kanan     tengah   trapesium    Simpson
  10  4.833e-02  5.167e-02  8.333e-04  1.667e-03  1.110e-16
 100  4.983e-03  5.017e-03  8.333e-06  1.667e-05  0.000e+00
1000  4.998e-04  5.002e-04  8.333e-08  1.667e-07  0.000e+00

quad exp(-x^2) = 0.746824132812427   galat perkiraan 8.3e-15
acuan erf      = 0.746824132812427   selisih +1.1e-16
quad x^2       = 0.33333333333333337

galat |Simpson - acuan| untuk exp(-x^2)
  n=  2  3.563e-04
  n=  4  3.125e-05
  n=  8  1.988e-06
  n= 16  1.246e-07
  n= 32  7.795e-09

galat |trapesium - acuan| untuk exp(-x^2)
  n=  10  6.133e-04
  n= 100  6.131e-06
  n=1000  6.131e-08
\end{verbatim}

Empat hal terbaca. Selisih terbagi dua setiap kali $n$ berlipat dua, sesuai $1/1024$
pada baris terakhir. Galat kiri dan kanan turun sepuluh kali lipat, titik tengah dan
trapesium seratus kali lipat. Kolom Simpson berisi $1.110\times10^{-16}$ lalu nol, sisa
pembulatan, dan galat Simpson untuk $e^{-x^2}$ turun sekitar enam belas kali lipat
setiap $n$ berlipat dua.

Kode itu tidak membuktikan satu pun teorema. Ia memeriksa bahwa rumus bab ini memberi
angka yang disebutkan. Bukti keterintegralan $x^2$ ada pada perhitungan $U_n-L_n=1/n$
di bagian sebelumnya: keluaran program adalah data, bukan argumen.

\begin{untukdosen}
Urutan mengajar yang saya sarankan: mulai dari dua jawaban yang berbeda di pembuka, lalu
definisi partisi dan hitung $U$, $L$ untuk $x^2$ dengan empat potongan di papan. Setelah
itu lemma penghalusan, kriteria, baru kelas fungsi. Jangan menunda teorema dasar
kalkulus: mahasiswa perlu melihat kerja keras kriteria berbuah menjadi alat hitung.

Dua miskonsepsi yang sering muncul. Pertama, integral dianggap sama dengan antiturunan,
sehingga $F'=f$ dianggap berlaku karena definisi, padahal itu teorema. Kedua, fungsi
dengan banyak titik diskontinu dianggap pasti tidak terintegralkan, karena contoh
Dirichlet terlalu sering dipakai sebagai satu-satunya contoh. Fungsi Thomae menghapus
anggapan itu dalam satu halaman.

Aktivitas papan tulis: minta mahasiswa menebak $\int_0^1x^2dx$ hanya dari $U_n$ dan
$L_n$ untuk $n=1,2,4$, lalu menyebutkan berapa $n$ yang diperlukan agar galat di bawah
$10^{-3}$ tanpa menghitung integralnya. Jawaban $n>1000$ muncul dari $1/n$, dan momen
itu menjelaskan mengapa antiturunan lebih berharga daripada kuadratur. Tanda siap
lanjut: mahasiswa dapat menjelaskan perbedaan arah kedua bagian teorema dasar kalkulus,
dan menyebutkan satu fungsi terintegralkan yang tidak kontinu.
\end{untukdosen}

\section{Ringkasan}
\tingkat{dasar}

\begin{ringkasan}
\begin{itemize}
\item Integral Riemann adalah satu-satunya bilangan yang terkurung di antara semua
      jumlah bawah dan semua jumlah atas; bila kedua kurungan bertemu, nilainya disebut
      integral.
\item Kriteria keterintegralan menggantikan pemeriksaan tak berhingga banyak partisi
      dengan satu pertidaksamaan $U(\mathcal{P},f)-L(\mathcal{P},f)<\varepsilon$.
\item Fungsi kontinu dan fungsi monoton terintegralkan; fungsi Dirichlet tidak, karena
      setiap potongannya memuat rasional dan irasional, sedangkan fungsi Thomae
      terintegralkan walaupun terputus di setiap bilangan rasional.
\item Fungsi akumulasi $F(x)=\int_a^xf$ selalu kontinu dan turunannya $f$ di setiap
      titik tempat $f$ kontinu, jadi setiap fungsi kontinu punya antiturunan.
\item Bagian kedua menyatakan $\int_a^bF'=F(b)-F(a)$; buktinya teorema nilai rata-rata
      ditambah teleskopik, dengan syarat $F'$ terintegralkan.
\item Substitusi dan integral parsial adalah aturan rantai dan aturan perkalian yang
      dibaca terbalik. Integral tak wajar adalah limit integral Riemann, dan uji
      integral menghubungkannya dengan deret Bab 5.
\end{itemize}
\end{ringkasan}

\section{Latihan}
\tingkat{menengah}

\begin{latihan}
\begin{enumerate}
  \item Hitung $U_n$ dan $L_n$ untuk $f(x)=x$ pada $[0,1]$ dengan partisi seragam,
        nyatakan keduanya sebagai pecahan dalam $n$, lalu tentukan $U_n-L_n$. Simpulkan
        bahwa $f$ terintegralkan dan hitung integralnya.
  \item Buktikan bahwa $f$ terintegralkan pada $[a,b]$ berakibat $|f|$ terintegralkan
        dan $\bigl|\int_a^bf\bigr|\le\int_a^b|f|$. Tunjukkan kebalikannya salah dengan
        satu contoh.
  \item Nilai klaim ini: ``Setiap fungsi terbatas pada $[a,b]$ terintegralkan, sebab
        jumlah atas dan jumlah bawahnya selalu berhingga.'' Benar, salah, atau belum
        cukup bukti? Sebutkan bagian mana yang didukung bukti dan bagian mana yang
        tidak.
  \item Buktikan bahwa fungsi naik pada $[a,b]$ terintegralkan dengan menunjukkan
        $U(\mathcal{P},f)-L(\mathcal{P},f)\le(b-a)\bigl(f(b)-f(a)\bigr)/n$ untuk
        partisi seragam $n$ bagian. Jelaskan mengapa batas itu tidak berlaku untuk
        fungsi yang tidak monoton.
  \item Hitung $\int_0^1(3x^2-2x+1)\,dx$ dengan teorema dasar kalkulus, lalu periksa
        hasilnya memakai aturan Simpson dengan $n=100$ dan fungsi
        \texttt{scipy.integrate.quad}. Laporkan ketiga angka dan selisihnya.
  \item Hitung $\int_1^{10}x^{-2}dx$ dan $\int_1^{\infty}x^{-2}dx$. Laporkan selisih
        keduanya, lalu jelaskan mengapa selisih itu sama dengan $1/10$ tanpa
        menghitung ulang.
  \item Buktikan bahwa bila $f$ kontinu, $f\ge0$ pada $[a,b]$, dan $\int_a^bf=0$, maka
        $f(x)=0$ untuk setiap $x\in[a,b]$. Tunjukkan juga bahwa pernyataan itu gagal
        bila $f$ hanya terintegralkan.
  \item Hitung $\int_0^{\pi/2}\sin x\cos x\,dx$ dengan dua cara: substitusi
        $u=\sin x$, dan identitas $\sin2x=2\sin x\cos x$. Bandingkan hasilnya dengan
        nilai numerik dari \texttt{quad}.
  \item Tunjukkan bahwa fungsi Dirichlet pada $[0,1]$ memenuhi
        $U(\mathcal{P},D)-L(\mathcal{P},D)=1$ untuk setiap partisi, lalu jelaskan
        mengapa fakta itu cukup untuk menyimpulkan ketidakterintegralannya.
  \item Taksir $\sum_{n=101}^{\infty}1/n^2$ dengan batas integral dan bandingkan dengan
        nilai sebenarnya $9.950167\times10^{-3}$; berapa galat relatifnya?
\end{enumerate}

\noindent\textbf{Petunjuk.} Butir 3: klaim memuat dua pernyataan berbeda, tentang
keterbatasan jumlah dan tentang kesamaan kedua integral. Butir 7: gunakan
kekontinuan di satu titik untuk membuat integral kecil pada selang sempit di sekitar
titik itu. Butir 10: pakai $\sum_{n=N+1}^\infty1/n^2\le\int_N^\infty x^{-2}dx=1/N$
dengan $N=100$.
\end{latihan}

\begin{untukpeneliti}
Yang benar-benar terbukti di bab ini: kriteria keterintegralan dengan kedua arahnya,
keterintegralan fungsi kontinu (memakai kekontinuan seragam dari Bab 11), fungsi
monoton, fungsi Thomae, kedua bagian teorema dasar kalkulus, teorema nilai rata-rata
untuk integral, dan status substitusi serta integral parsial sebagai akibat, bukan
aksioma.

Yang dikutip tanpa bukti dan sebaiknya disebut begitu di kelas: kriteria Lebesgue
(Rudin, Teorema 11.33(b); Abbott, Bab 7), contoh Volterra tentang turunan terbatas yang
tidak terintegralkan, dan ketiadaan antiturunan elementer untuk $e^{-x^2}$ (teorema
Liouville).

Sumber primer: Rudin, \emph{Principles of Mathematical Analysis}, Bab 6, bagian
Definition and Existence of the Integral (Teorema 6.4, 6.6, 6.8, 6.9, 6.10) dan bagian
Integration and Differentiation (6.19--6.22). Rudin menyusunnya untuk integral
Riemann-Stieltjes; integral Riemann adalah kasus $\alpha(x)=x$. Abbott,
\emph{Understanding Analysis}, Bab 7, lebih lambat tetapi lebih terbaca; Stewart hanya
untuk aturan hitung.

Pertanyaan terbuka: seberapa tajam batas galat kuadratur bila hanya $|f''|$ yang
diketahui, dan bagaimana memperkirakan galat tanpa turunan sama sekali; jawaban
standarnya kuadratur adaptif dan selisih dua tingkat penghalusan. Batas keberlakuan:
fungsi terbatas pada selang tertutup terbatas, integral tak wajar sebagai limit, tanpa
teori ukur, dan integral bernilai vektor tidak disinggung.
\end{untukpeneliti}

\peta{Bab 13 pada Rudin, Bab 6: bagian Definition and Existence of the Integral
(partisi, jumlah atas-bawah, Lemma penghalusan sebagai Teorema 6.4, kriteria sebagai
6.6, fungsi kontinu dan monoton sebagai 6.8 dan 6.9) dan bagian Integration and
Differentiation (6.19 perubahan variabel, 6.21 teorema dasar kalkulus, 6.22 integral
parsial). Pada Abbott, Bab 7, dari The Definition of the Riemann Integral sampai
Lebesgue's Criterion. Pada Stewart: The Substitution Rule dan Average Value of a
Function (Bab 5), Integration by Parts dan Improper Integrals (Bab 7), serta 11.3
The Integral Test, jembatan ke Bab 5 buku ini.}